\chapter{绪论}

\section{研究背景与意义}

本节用于说明课题产生的背景、实际问题和研究价值。写作时应从具体应用场景进入，避免只堆砌宏观口号。可以先描述业务痛点，再说明为什么需要设计并实现一个系统来解决该问题。

\section{国内外研究现状}

本节用于梳理与课题相关的研究和工程实践。建议按照研究方向或技术路线组织内容，并在正文中使用参考文献引用，例如：相关研究为本文的系统设计提供了参考\cite{ref1}。

\section{研究内容与论文结构}

本节用于说明论文主要工作和章节安排。正式论文中可简要列出系统需求分析、总体设计、详细实现、测试分析和总结展望等内容。

\chapter{系统需求分析}

\section{需求概述}

本章用于说明系统面向的用户角色、业务流程和功能边界。需求分析应聚焦真实问题，明确系统需要解决什么、不解决什么。

\section{功能需求}

\begin{enumerate}
\item 用户管理与权限控制。
\item 核心业务数据维护。
\item 业务流程处理与状态流转。
\item 统计分析、查询检索或智能辅助功能。
\end{enumerate}

\section{非功能需求}

非功能需求可从性能、安全性、可靠性、可维护性和易用性等方面展开。若论文中给出指标，应保证指标可测试、可解释。

\begin{figure}[H]
\centering
\IfFileExists{figures/generated/fig-2-1.pdf}{\includegraphics[width=0.82\textwidth]{figures/generated/fig-2-1.pdf}}{%
\begin{minipage}{0.9\textwidth}
\lstinputlisting[style=mermaid]{figures/mermaid/fig-2-1.mmd}
\end{minipage}}
\caption{系统典型业务流程}\label{fig:2-1}
\end{figure}

\chapter{系统总体设计}

\section{总体架构设计}

本节用于说明系统整体架构、模块边界和数据流向。图 3-1 给出一个通用分层架构示例，正式使用时应替换为真实系统架构。

\begin{figure}[H]
\centering
\IfFileExists{figures/generated/fig-3-1.pdf}{\includegraphics[width=0.82\textwidth]{figures/generated/fig-3-1.pdf}}{%
\begin{minipage}{0.9\textwidth}
\lstinputlisting[style=mermaid]{figures/mermaid/fig-3-1.mmd}
\end{minipage}}
\caption{系统总体架构}\label{fig:3-1}
\end{figure}

\section{数据库设计}

数据库设计应结合实体关系、字段含义、索引约束和状态流转展开。表 3-1 给出通用示例，正式论文中应替换为真实表结构。

\begingroup
\zihao{5}\setstretch{1.15}\setlength{\tabcolsep}{3pt}\renewcommand{\arraystretch}{1.25}
\begin{longtable}{|>{\RaggedRight\arraybackslash}p{0.18\textwidth}|>{\RaggedRight\arraybackslash}p{0.25\textwidth}|>{\RaggedRight\arraybackslash}p{0.39\textwidth}|}
\caption{核心数据表示例}\label{tab:3-1}\\
\hline
\textbf{表名} & \textbf{主要字段} & \textbf{说明} \\ \hline
\endfirsthead
\multicolumn{3}{c}{续表~\thetable}\\
\hline
\textbf{表名} & \textbf{主要字段} & \textbf{说明} \\ \hline
\endhead
用户表 & \texttt{id}、\texttt{name}、\texttt{role} & 保存用户基础信息和角色信息 \\ \hline
业务记录表 & \texttt{id}、\texttt{status}、\texttt{created\_at} & 保存核心业务记录及其状态 \\ \hline
操作日志表 & \texttt{id}、\texttt{operator\_id}、\texttt{content} & 保存关键操作日志，便于审计和追溯 \\ \hline
\end{longtable}
\endgroup

\chapter{系统详细设计与实现}

\section{核心模块实现}

本章应是全文重点。建议每个主要功能模块包含模块说明、页面或流程截图、关键代码片段和实现难点。下面给出截图占位示例，放入同名图片后会自动替换。

\begin{figure}[H]
\centering
\renderscreenshot{4-1}{screenshot-module-page.png}{核心功能页面截图占位}{请填写页面或功能来源}
\end{figure}

\section{关键代码示例}

代码片段应服务于论文说明，不宜直接粘贴大段源码。示例如下：

\begin{lstlisting}[style=thesiscode,language=python]
def validate_status(status: str) -> bool:
    allowed = {"draft", "processing", "finished"}
    return status in allowed
\end{lstlisting}

\chapter{系统测试与结果分析}

\section{测试环境}

本节用于说明硬件环境、软件环境、数据库版本、浏览器版本和测试数据来源。正式写作时应使用真实可复核的信息。

\section{功能测试}

\begingroup
\zihao{5}\setstretch{1.15}\setlength{\tabcolsep}{3pt}\renewcommand{\arraystretch}{1.25}
\begin{longtable}{|>{\RaggedRight\arraybackslash}p{0.18\textwidth}|>{\RaggedRight\arraybackslash}p{0.33\textwidth}|>{\RaggedRight\arraybackslash}p{0.28\textwidth}|}
\caption{功能测试用例示例}\label{tab:5-1}\\
\hline
\textbf{编号} & \textbf{测试内容} & \textbf{预期结果} \\ \hline
\endfirsthead
\multicolumn{3}{c}{续表~\thetable}\\
\hline
\textbf{编号} & \textbf{测试内容} & \textbf{预期结果} \\ \hline
\endhead
T01 & 登录系统并访问首页 & 页面正常显示 \\ \hline
T02 & 新增一条核心业务记录 & 记录保存成功且状态正确 \\ \hline
T03 & 查询历史业务记录 & 返回符合条件的数据 \\ \hline
\end{longtable}
\endgroup

\chapter{总结与展望}

\section{工作总结}

本节用于总结论文完成的主要工作、系统实现结果和测试结论。请避免夸大系统成熟度，应与实际实现保持一致。

\section{未来展望}

本节用于说明系统存在的不足和后续改进方向。可以从功能完善、性能优化、部署推广和数据安全等角度展开。
